\section*{الفصل \ref{ch:g11:functional-groups} --- المجموعات الوظيفية والعائلات}

\begin{solution}{exo:g11:functional-groups:1}
هيدروكسيل، \omterm{def:g11:functional-groups:alcohol}{كحول}؛ كربونيل داخل السلسلة، \omterm{def:g11:functional-groups:carbonyl}{كيتون}؛ كربوكسيل، \omterm{def:g11:functional-groups:carboxylic-acid}{حمض
كربوكسيلي}؛ مجموعة \omterm{def:g11:functional-groups:carboxylic-acid}{إستر}، \omterm{def:g11:functional-groups:carboxylic-acid}{إستر}؛ مجموعة أمينو \ce{-NH2}، \omterm{def:g11:functional-groups:amine}{أمين}.
\end{solution}

\begin{solution}{exo:g11:functional-groups:2}
الميثانول؛ بروبان-1-ول؛ حمض الميثانويك؛ البروبانال؛ البروبانون.
\end{solution}

\begin{solution}{exo:g11:functional-groups:3}
\omterm{def:g11:functional-groups:alcohol}{كحول}؛ \omterm{def:g11:functional-groups:carbonyl}{ألدهيد}؛ \omterm{def:g11:functional-groups:carboxylic-acid}{إستر}؛ \omterm{def:g11:functional-groups:amine}{أميد}؛ \omterm{def:g11:functional-groups:halogenoalkane}{ألكان مهلجن}؛ \omterm{def:g11:functional-groups:halogenoalkane}{ألكين}.
\end{solution}

\begin{solution}{exo:g11:functional-groups:4}
البروبانال \ce{CH3-CH2-CHO}، والبروبانون \ce{CH3-CO-CH3}. والبروبانال هو
\omterm{def:g11:functional-groups:carbonyl}{الألدهيد}: فكربون \ce{C=O} فيه في طرف السلسلة، مرتبط \omterm{def:g7:atoms-and-molecules:atom}{بذرة}
هيدروجين؛ وفي البروبانون يقع بين ذرتي كربون.
\end{solution}

\begin{solution}{exo:g11:functional-groups:5}
أولي (الكربون الحامل للمجموعة \ce{-OH} له جار كربوني واحد)؛
ثانوي (جاران)؛ ثالثي (ثلاثة).
\end{solution}

\begin{solution}{exo:g11:functional-groups:6}
\setchemfig{atom sep=1.5em}
بنتان-3-ول \chemfig{-[:30]-[:-30](-[:-90]OH)-[:30]-[:-30]}؛
حمض 2-ميثيل البوتانويك \chemfig{-[:30]-[:-30](-[:-90])-[:30](=[:90]O)-[:-30]OH}؛
هكسان-2-ون \chemfig{-[:30](=[:90]O)-[:-30]-[:30]-[:-30]-[:30]}؛
إيثانوات البروبيل \chemfig{-[:30](=[:90]O)-[:-30]O-[:30]-[:-30]-[:30]}.
\end{solution}

\begin{solution}{exo:g11:functional-groups:7}
\ce{C3H6O}. نعم، هما متماكبان، لكن بمجموعتين مختلفتين:
\omterm{def:g11:functional-groups:carbonyl}{ألدهيد} \omterm{def:g11:functional-groups:carbonyl}{وكيتون} (وكلتاهما مجموعة كربونيل، في موضعين مختلفين).
\end{solution}

\begin{solution}{exo:g11:functional-groups:8}
\ce{HCOO-CH2-CH3}، ميثانوات الإيثيل؛ \ce{CH3-COO-CH2-CH2-CH3}، إيثانوات
البروبيل.
\end{solution}

\begin{solution}{exo:g11:functional-groups:9}
إيثانوات الميثيل (حمض الإيثانويك، والميثانول)؛ وبروبانوات الإيثيل (حمض
البروبانويك، والإيثانول)؛ وميثانوات الميثيل (حمض الميثانويك، والميثانول).
\end{solution}

\begin{solution}{exo:g11:functional-groups:10}
\setchemfig{atom sep=1.5em}
\begin{center}
\small
\begin{tabular}{cccc}
\chemfig{-[:30]-[:-30]-[:30]-[:-30]OH} &
\chemfig{-[:30](-[:90]OH)-[:-30]-[:30]} &
\chemfig{-[:30](-[:90])-[:-30]-[:30]OH} &
\chemfig{-[:0](-[:90]OH)(-[:-90])-[:0]} \\
بوتان-1-ول & بوتان-2-ول & 2-ميثيل بروبان-1-ول & 2-ميثيل بروبان-2-ول \\
أولي & ثانوي & أولي & ثالثي
\end{tabular}
\end{center}
\end{solution}

\begin{solution}{exo:g11:functional-groups:11}
\omterm{def:g11:functional-groups:amine}{أميد} (الإيثان \omterm{def:g11:functional-groups:amine}{أميد}). \emph{-ين}. الأحماض الكربوكسيلية (\ce{C=O}
و\ce{-OH} على الكربون نفسه)، \omterm{def:g11:functional-groups:carboxylic-acid}{والإسترات} (\ce{C=O} و\ce{-O-}) والأميدات
(\ce{C=O} و\ce{N}): ثلاث عائلات.
\end{solution}

\begin{solution}{exo:g11:functional-groups:12}
الأسبرين: مجموعة كربوكسيل (\omterm{def:g11:functional-groups:carboxylic-acid}{حمض كربوكسيلي}) ومجموعة \omterm{def:g11:functional-groups:carboxylic-acid}{إستر}.
الباراسيتامول: مجموعة فينول \ce{-OH} على الحلقة ومجموعة \omterm{def:g11:functional-groups:amine}{أميد}
\ce{NH-CO}.
\end{solution}

\begin{solution}{exo:g11:functional-groups:13}
\ce{C2H6O}. الإيثانول وحده، بمجموعته \ce{O-H}، يكوّن روابط هيدروجينية
بين جزيئاته (فليس على أكسجين الميثوكسي ميثان هيدروجين).
والإيثانول يغلي عند درجة أعلى، \qty{78}{\celsius}؛ والميثوكسي ميثان عند
\qty{-25}{\celsius}.
\end{solution}

\begin{solution}{exo:g11:functional-groups:14}
مجموعة هيدروكسيل ومجموعة كربوكسيل. للسلسلة ثلاث \omterm{def:g7:atoms-and-molecules:atom}{ذرات} كربون، كربون الحمض هو
الكربون 1 والمجموعة \ce{-OH} على الكربون 2:
حمض 2-هيدروكسي البروبانويك. والغليسين: حمض 2-أمينو الإيثانويك (ويُكتب كثيرًا
حمض أمينو الإيثانويك).
\end{solution}

\begin{solution}{exo:g11:functional-groups:15}
الإيثانول \ce{CH3-CH2-OH}، \ce{C2H6O}، \omterm{def:g11:functional-groups:alcohol}{كحول}؛ والإيثانال \ce{CH3-CHO}،
\ce{C2H4O}، \omterm{def:g11:functional-groups:carbonyl}{ألدهيد}؛ وحمض الإيثانويك \ce{CH3-COOH}، \ce{C2H4O2}،
\omterm{def:g11:functional-groups:carboxylic-acid}{حمض كربوكسيلي}. فمن الإيثانول إلى الإيثانال تُفقد ذرتا هيدروجين؛
ومن الإيثانال إلى حمض الإيثانويك تُكتسب \omterm{def:g7:atoms-and-molecules:atom}{ذرة} أكسجين واحدة.
\end{solution}

\begin{solution}{pb:g11:functional-groups:1}
\textbf{1.} A: مجموعة \omterm{def:g11:functional-groups:carboxylic-acid}{إستر}؛ B: هيدروكسيل؛ C: كربوكسيل؛ D: كربونيل في
طرف السلسلة؛ E: مجموعة \omterm{def:g11:functional-groups:carboxylic-acid}{إستر}.

\textbf{2.} A \omterm{def:g11:functional-groups:carboxylic-acid}{إستر}، وB \omterm{def:g11:functional-groups:alcohol}{كحول}، وC \omterm{def:g11:functional-groups:carboxylic-acid}{حمض كربوكسيلي}، وD \omterm{def:g11:functional-groups:carbonyl}{ألدهيد}، وE \omterm{def:g11:functional-groups:carboxylic-acid}{إستر}.

\textbf{3.} A (الموز) وE (الأناناس): روائح فاكهية.

\textbf{4.} أولي: فالكربون الحامل للمجموعة \ce{-OH} مرتبط \omterm{def:g7:atoms-and-molecules:atom}{بذرة} كربون واحدة.

\textbf{5.} \ce{C6H12O}؛ مثلًا هكسان-2-ون،
\ce{CH3-CO-CH2-CH2-CH2-CH3}.

\textbf{6.} السلسلة الحاملة للمجموعة \ce{C-OH} فيها أربع \omterm{def:g7:atoms-and-molecules:atom}{ذرات} كربون، مرقّمة
من جهة \ce{-OH}، وعلى الكربون 3 مجموعة ميثيل: 3-ميثيل بوتان-1-ول.

\textbf{7.} C: حمض الإيثانويك؛ D: الهكسانال.

\textbf{8.} بوتانوات الإيثيل، المرتبط بحمض البوتانويك والإيثانول.

\textbf{9.} \emph{إيثانوات}: الجزء الحمضي \ce{CH3-CO}، من حمض
الإيثانويك؛ و\emph{3-ميثيل البوتيل}: الجزء الكحولي، من 3-ميثيل بوتان-1-ول.

\textbf{10.} C (حمض الإيثانويك) وB (3-ميثيل بوتان-1-ول).

\textbf{11.} A \ce{C7H14O2}؛ B \ce{C5H12O}؛ C \ce{C2H4O2}.

\textbf{12.} $\ce{C2H4O2} + \ce{C5H12O}$ يحوي \ce{C7H16O3}، أي
\ce{H2O} واحدًا زيادة على \ce{C7H14O2}: إنه الماء.

\textbf{13.} \ce{C2H4O2 + C5H12O -> C7H14O2 + H2O}: C 7 = 7، H 16 = 16،
O 3 = 3.

\textbf{14.} $M(\text{B}) = 5 \times 12.0 + 12 \times 1.0 + 16.0 =
\qty{88.0}{g/mol}$؛ $M(\text{C}) = \qty{60.0}{g/mol}$.

\textbf{15.} $60.0 + 88.0 - 18.0 = \qty{130.0}{g/mol}$.

\textbf{16.} $7 \times 12.0 + 14 \times 1.0 + 2 \times 16.0 =
\qty{130.0}{g/mol}$: القيمة نفسها.

\textbf{17.} E، \ce{C6H12O2}: $6 \times 12.0 + 12 \times 1.0 + 2 \times
16.0 = \qty{116.0}{g/mol}$.

\textbf{18.} \ce{C7H14O2}: الصيغة نفسها التي للمركب A، فهو متماكب له، وله
\omterm{def:g10:the-mole:molar-mass}{الكتلة المولية} نفسها. \omterm{def:g10:the-mole:molar-mass}{فالكتلة المولية} وحدها لا تميّز بين \omterm{def:g11:organic-skeletons:isomer}{المتماكبات}.

\textbf{19.} للمركب A ولحمض الهبتانويك \omterm{def:g7:atoms-and-molecules:atom}{الذرات} نفسها لكن مجموعتيهما
مختلفتان، \omterm{def:g11:functional-groups:carboxylic-acid}{إستر} \omterm{def:g11:functional-groups:carboxylic-acid}{وحمض كربوكسيلي}: فأحدهما رائحته رائحة الموز، والآخر
زنخة. فالمجموعة، وموضعها، هما الحاسمان.

\textbf{20.} \qty{130.0}{g/mol}.
\end{solution}
